\documentclass[10pt]{ltjsarticle}
\usepackage{tem}
\usepackage{hyperref}
\usepackage{listings}
\pagestyle{fancy}
\lhead{Infinite Generation Two}
\rhead{Native Pi migration evidence}
\setlength{\headheight}{16pt}
\setlength{\headsep}{30pt}
\definecolor{RemovedBg}{HTML}{FCE8E6}
\definecolor{AddedBg}{HTML}{E6F4EA}
\definecolor{CodeKeyword}{HTML}{633194}
\definecolor{CodeString}{HTML}{075E54}
\definecolor{CodeComment}{HTML}{586A6A}
\lstdefinelanguage{TypeScript}{
  morekeywords={import,from,export,default,type,function,const,let,if,else,return,async,await,true,false,null,undefined,new,throw,void,string,boolean},
  sensitive=true,
  morecomment=[l]{//},
  morecomment=[s]{/*}{*/},
  morestring=[b]",
  morestring=[b]',
  morestring=[b]`
}
\lstset{language=TypeScript,
  basicstyle=\ttfamily\footnotesize\addfontfeatures{RawFeature=-calt},
  keywordstyle=\color{CodeKeyword}\bfseries,
  stringstyle=\color{CodeString},commentstyle=\color{CodeComment},
  breaklines=true,breakatwhitespace=false,columns=fullflexible,keepspaces=true,
  showstringspaces=false,numbers=left,numberstyle=\tiny\color{gray},
  numbersep=8pt,xleftmargin=18pt,xrightmargin=4pt,frame=single,
  rulecolor=\color{gray!35},tabsize=2,aboveskip=6pt,belowskip=6pt}
\begin{document}
\title{Native Pi Extension Migration\\\large{--- Code Change Evidence ---}}
\author{Fides Thérèse}
\date{2026-10-22}
\maketitle
\thispagestyle{plain}

\section{Change metadata}
Timestamp: \texttt{20261022T093200Z}. Delivery mode: Standard.
Triggers: CRUD read behavior (bundled prompt/profile exposure) and bug fixes
(host incompatibility, cancellation status, and mid-run toggle state).
Baseline commit: \path{3a701c334dd15e59c6ac663d93d313489385bbc5}.
The highlighted regions are extracted from the final staged source diff;
line numbers refer to the corresponding original or added source file.
The SHA-256 of the staged code/config/test diff is:\par\noindent
\texttt{22bd49efe8a29f4f25e8dcf486b7c131ddbdefbc8d81d896247d9937f926b490}

\section{Intent and affected surface}
Replace DeepSeek Harness-only registration with a native Pi extension while
preserving the original prompt and license exactly. Pi must retain its base
instructions, tools, project context, and approval mechanisms. The English
README replaces obsolete installation instructions and promotional material.

\begin{description}
\item[Adapter.] \path{index.ts}: default factory, named prompt section,
commands, startup flag, assistant-message observation, automatic compaction
failure outcomes, final settlement, and keyed UI cleanup.
\item[Helpers.] \path{src/profile.ts}: integrity loader, metadata builder,
TypeBox schemas, and profile-tool result. \path{src/armor.ts}: text extraction,
inherited bilingual score heuristic, and English status formatting.
\item[Removed host-specific surface.] Legacy entry point, browser React badge,
Cordis declaration, and four Harness install/uninstall scripts. The unused
sponsor image and old Harness note are removed; Git retains their history.
\item[Configuration and verification.] Pi package manifest, lockfile, strict
TypeScript configuration, exact-byte Git attributes, SDK/unit tests, and
CLI/package-install smoke tests. The prompt and MIT license have no diff.
\end{description}

\section{Highlighted code changes}
Pale red denotes removed code; pale green denotes added code. Keywords,
strings, and comments retain independent syntax colors. The original profile
used the Harness \texttt{output} renderer and returned a raw value; the native
tool explicitly returns Pi's content, details, and structured result. Routine
unchanged metadata and wholesale removed installers are not dumped here.
\subsection{Prompt registration}
\noindent\textbf{Removed --- \texttt{index.js}.}\par\smallskip
\begin{lstlisting}[backgroundcolor=\color{RemovedBg},firstnumber=128]
  ctx.effect(() => ctx.systemPrompt.section({
    name: "infinite-gen-2:global-system-prompt",
    order: 100,
    text: PROMPT_TEXT,
  }));
  ctx.effect(() => ctx.tools.register(profileTool));
\end{lstlisting}
\noindent\textbf{Added --- \texttt{index.ts}.}\par\smallskip
\begin{lstlisting}[backgroundcolor=\color{AddedBg},firstnumber=112]
  pi.on("before_agent_start", (event) => {
    if (state.enabled) {
      event.systemPromptOptions.sections[PROMPT_SECTION] = promptText;
    } else {
      delete event.systemPromptOptions.sections[PROMPT_SECTION];
    }
  });
\end{lstlisting}
\subsection{Tool result contract}
\noindent\textbf{Removed --- \texttt{index.js}.}\par\smallskip
\begin{lstlisting}[backgroundcolor=\color{RemovedBg},firstnumber=6]
const objectOutput = {
  schema: { type: "object", additionalProperties: true },
  render: (_args, value) => [{ type: "text", text: JSON.stringify(value) }],
};
\end{lstlisting}
\noindent\textbf{Added --- \texttt{src/profile.ts}.}\par\smallskip
\begin{lstlisting}[backgroundcolor=\color{AddedBg},firstnumber=95]
  return {
    name: "infinite_gen2_profile",
    label: "Infinite Gen 2 profile",
    description: "Return the unchanged bundled Infinite Generation Two system prompt and upstream calibration metadata.",
    parameters: PROFILE_PARAMETERS,
    outputSchema: PROFILE_SCHEMA,
    annotations: {
      readOnlyHint: true,
      destructiveHint: false,
      idempotentHint: true,
      openWorldHint: false,
    },
    /**
     * Return a fresh profile result for the bundled prompt.
     *
     * This operation has no arguments or external side effects and exposes the
     * same value to the model, UI, and programmatic tool callers.
     *
     * return result: Promise<AgentToolResult<InfiniteGen2Profile>>; complete profile result.
     */
    async execute() {
      const profile = buildProfile(promptText);
      return {
        content: [{ type: "text", text: JSON.stringify(profile) }],
        details: profile,
        structuredContent: profile,
      };
    },
  };
\end{lstlisting}
\subsection{Run and cancellation lifecycle}
\noindent\textbf{Removed --- \texttt{index.js}.}\par\smallskip
\begin{lstlisting}[backgroundcolor=\color{RemovedBg},firstnumber=110]
function armorProjectionApply(state, event) {
  if (!event || typeof event !== "object") return state;
  if (event.type === "user/message") {
    return { running: true, verdict: null, words: [] };
  }
  if (event.type === "assistant/message") {
    const text = armorTextOf(event?.data?.message?.content);
    if (!text.trim()) return state;
    const scored = armorScore(text);
    return { running: false, verdict: scored.verdict, words: scored.words };
  }
  return state;
}
\end{lstlisting}
\noindent\textbf{Added --- \texttt{index.ts}.}\par\smallskip
\begin{lstlisting}[backgroundcolor=\color{AddedBg},firstnumber=120]
  pi.on("agent_start", (_event, ctx) => {
    state = { ...state, running: true, score: null, outcome: null };
    publishArmorStatus(ctx);
  });
  pi.on("message_end", (event, ctx) => {
    if (!state.enabled || event.message.role !== "assistant") return;
    const message = event.message;
    if (message.stopReason === "error" || message.stopReason === "aborted") {
      state = { ...state, score: null, outcome: message.stopReason };
    } else {
      const text = extractResponseText(message.content);
      if (!text.trim()) return;
      state = { ...state, score: scoreResponseText(text), outcome: null };
    }
    publishArmorStatus(ctx);
  });
  pi.on("session_compact_failed", (event, ctx) => {
    if (!state.enabled || !state.running || event.reason === "manual") return;
    state = { ...state, score: null, outcome: event.aborted ? "aborted" : "error" };
    publishArmorStatus(ctx);
  });
  pi.on("agent_settled", (_event, ctx) => {
    state = { ...state, running: false };
    publishArmorStatus(ctx);
  });
\end{lstlisting}

\subsection{Mid-run enablement invariant}
The added command transition below preserves the actual run's busy state.
Enablement affects scoring and the next prompt, not whether Pi has settled.
\noindent\textbf{Added --- \texttt{index.ts}.}\par\smallskip
\begin{lstlisting}[backgroundcolor=\color{AddedBg}]
      if (action === "on" || action === "off") {
        state = { ...state, enabled: action === "on", score: null, outcome: null };
        publishArmorStatus(ctx);
        message = formatArmorStatus(state);
\end{lstlisting}

\section{Reasoning}
A stable default factory and explicit package resource manifest replace the
unavailable Harness/Cordis/browser APIs. Assigning the same owned section is
idempotent and preserves all unrelated instructions; disabling deletes only
that section. The integrity loader hashes raw bytes before UTF-8 decoding,
and the Git attribute prevents checkout-time newline conversion.

The native tool's TypeBox input rejects extra arguments. Its output schema,
model-facing text, details, and structured JSON describe the same fresh object.
It is read-only and remains callable when prompt injection is disabled.

The adapter records busy state until \texttt{agent\_settled}, even across
retries, queued work, and toggle commands. A failed or cancelled automatic
compaction can occur after a successful assistant message without another
assistant response; its event therefore clears stale pass status. Manual or
idle compaction is ignored. Errors and cancellations never become a passing
verdict. The detector observes text only and never rewrites output or
initiates a model retry.

\section{Verification}
Executed on Windows with Node.js 24.11.1 and the pinned Pi 0.99.1 development
API. All commands below completed successfully:
\begin{itemize}
\item \texttt{npm ci --ignore-scripts --no-audit --no-fund}: installed the
locked development dependencies. npm reported one upstream
\texttt{node-domexception} deprecation warning; no install scripts ran.
\item \texttt{npm test}: 31 SDK/unit tests passed, including real Pi loader,
runner, command, and tool execution; prompt tampering/missing assets;
section preservation; independent sessions; cancellation; and mid-run toggles.
\item \texttt{npm run check}: strict TypeScript check passed.
\item \texttt{npm run smoke}: three tests passed for local installation,
configured package discovery, directory loading, and explicit-file loading.
All CLI state was isolated from the user's agent directory.
\item \texttt{npm pack --dry-run}: seven required files, including the prompt
and both source helpers; no browser, installer, or sponsor assets.
\item Pi 1.4.0 was invoked through \texttt{npm exec} with offline resource
flags and this package as its explicit extension. Its help output included
\texttt{--infinite-gen2-off}; extension loading succeeded.
\item Node assertions compared the prompt and license bytes directly with
\texttt{git show HEAD:<path>}; both matched. The README passed English-only
and removed-advertisement checks. \texttt{git diff --cached --check} passed.
\end{itemize}
The first regression demonstrated that the legacy entry point had no Pi
default factory. Review regressions reproduced both a stale passing score
after automatic compaction cancellation and premature idle state on enablement;
the focused five-case suite passed after correction.

The prompt's frozen SHA-256 remains:\par\noindent
\texttt{02679eb726507d5a6bfe0a327c8469965440be7a7c5104cab57b8d1e529ec466}

\section{Residual limits}
No live DeepSeek or other paid model request was made. SDK lifecycle coverage
targets Pi 0.99.1; Pi 1.4.0 has a CLI-load smoke check, not the same complete
SDK suite. The inherited phrase detector can match quotations and does not
measure prompt effectiveness. Another extension can supersede structured
sections by replacing the full prompt. Toggles do not retract already-sent
instructions or prompt content returned by the profile tool.
\end{document}
